\documentclass[aps,preprint]{revtex4}%
\usepackage{amsfonts}
\usepackage{amsmath}
\usepackage{amssymb}
\usepackage[dvips]{graphicx}%
\setcounter{MaxMatrixCols}{30}
%TCIDATA{OutputFilter=latex2.dll}
%TCIDATA{Version=4.00.0.2321}
%TCIDATA{CSTFile=revtex4.cst}
%TCIDATA{Created=Tuesday, May 14, 2002 12:13:47}
%TCIDATA{LastRevised=Tuesday, May 14, 2002 14:47:33}
%TCIDATA{<META NAME="GraphicsSave" CONTENT="32">}
%TCIDATA{<META NAME="DocumentShell" CONTENT="Articles\SW\REVTeX 4">}
\newtheorem{theorem}{Theorem}
\newtheorem{acknowledgement}[theorem]{Acknowledgement}
\newtheorem{algorithm}[theorem]{Algorithm}
\newtheorem{axiom}[theorem]{Axiom}
\newtheorem{claim}[theorem]{Claim}
\newtheorem{conclusion}[theorem]{Conclusion}
\newtheorem{condition}[theorem]{Condition}
\newtheorem{conjecture}[theorem]{Conjecture}
\newtheorem{corollary}[theorem]{Corollary}
\newtheorem{criterion}[theorem]{Criterion}
\newtheorem{definition}[theorem]{Definition}
\newtheorem{example}[theorem]{Example}
\newtheorem{exercise}[theorem]{Exercise}
\newtheorem{lemma}[theorem]{Lemma}
\newtheorem{notation}[theorem]{Notation}
\newtheorem{problem}[theorem]{Problem}
\newtheorem{proposition}[theorem]{Proposition}
\newtheorem{remark}[theorem]{Remark}
\newtheorem{solution}[theorem]{Solution}
\newtheorem{summary}[theorem]{Summary}
\newenvironment{proof}[1][Proof]{\noindent\textbf{#1.} }{\ \rule{0.5em}{0.5em}}
\begin{document}
\preprint{ }
\title[Summary]{Summary of Content of Proposed Papers}
\author{Brian Weiner}
\affiliation{}
\keywords{}
\date{05/14/02}
\received{}

\startpage{1 }
\maketitle
\tableofcontents


\begin{enumerate}
\item \emph{All} $N$-particle states, $\breve{D}^{N}$, (pure and ensemble) can
be obtained from a reference positive operator $\breve{D}_{ref}\in
\mathcal{B}_{1+}\left(  \mathcal{H}^{N}\right)  $, by
\begin{equation}
\breve{D}^{N}=\hat{T}\left(  D_{ref}\right)  =\breve{T}\breve{D}_{ref}%
\breve{T}^{\dagger};\quad D_{ref}\left(  \mathsf{z}^{N},\mathsf{z}^{\prime
N}\right)  \overset{a.e.}{\neq}0\text{ and }rank\left\{  \breve{D}%
_{ref}\right\}  =\infty
\end{equation}
where $\left\{  \breve{T}\right\}  $ are diagonal operators
\begin{equation}
\breve{T}=\int T\left(  \mathsf{z}^{N}\right)  \left\vert \mathsf{z}%
^{N}\right\rangle \left\langle \mathsf{z}^{N}\right\vert d\mathsf{z}^{N}%
=\int\nu\left(  \mathsf{z}^{N}\right)  e^{i\omega\left(  \mathsf{z}%
^{N}\right)  }\left\vert \mathsf{z}^{N}\right\rangle \left\langle
\mathsf{z}^{N}\right\vert d\mathsf{z}^{N};\,\nu\left(  \mathsf{z}^{N}\right)
\geq0
\end{equation}
and $T\left(  \mathsf{z}^{N}\right)  \in L_{\infty}\left(  \mathsf{Z}%
^{N}\right)  $ and $T\left(  \mathsf{z}^{N}\right)  $ are symmetric functions
of the coordinates that refer to indistinguishable particles. This set of
transformations forms an abelian semi group, $\mathfrak{S}$, that is also a
$w^{\ast}$-algebra ($L_{\infty}\left(  \mathsf{Z}^{N}\right)  $ is a
representation of this algebra), which contains a group, $\mathfrak{G}_{pos}$,
of invertible elements
\begin{equation}
\mathfrak{G}_{pos}=\left\{  \breve{T}|T\left(  \mathsf{z}^{N}\right)
\overset{.}{\overset{a.e.}{\neq}0\Leftrightarrow\nu\left(  \mathsf{z}%
^{N}\right)  \overset{a.e.}{\neq}0}\right\}
\end{equation}
The group $\mathfrak{G}_{pos}$ is not sufficient to generate all states from a reference.

The diagonal operators can be expressed in second quantized notation as
\begin{equation}
\breve{T}=\int T\left(  \mathsf{z}^{N}\right)  \prod_{1\leq i<j\leq N}%
\delta^{C}\left(  \mathsf{z}_{i}-\mathsf{z}_{j}\right)  \Phi\left(
\mathsf{z}_{1}\right)  ^{\dagger}\Phi\left(  \mathsf{z}_{1}\right)  \cdots
\Phi\left(  \mathsf{z}_{N}\right)  ^{\dagger}\Phi\left(  \mathsf{z}%
_{N}\right)  d\mathsf{z}^{N}%
\end{equation}
where
\begin{equation}
\delta^{C}\left(  \mathsf{z}_{i}-\mathsf{z}_{j}\right)  =1-\delta\left(
\mathsf{z}_{i}-\mathsf{z}_{j}\right)
\end{equation}


\item The contraction maps
\begin{subequations}
\begin{align}
\Xi &  :\mathcal{B}_{1}\left(  \mathcal{H}^{N}\right)  \rightarrow
L_{1}\left(  \mathsf{Z}\right)  \subset\mathcal{B}_{1}\left(  \mathcal{H}%
^{1}\right) \\
L_{N}^{1}  &  :\mathcal{B}_{1}\left(  \mathcal{H}^{N}\right)  \rightarrow
\mathcal{B}_{1}\left(  \mathcal{H}^{1}\right) \\
L_{N}^{2}  &  :\mathcal{B}_{1}\left(  \mathcal{H}^{N}\right)  \rightarrow
\mathcal{B}_{1}\left(  \mathcal{H}^{2}\right)
\end{align}
are linear and thus their kernels are linear subspaces of $\mathcal{B}%
_{1}\left(  \mathcal{H}^{N}\right)  $. The maps $L_{N}^{1}$ and $L_{N}^{2}$
are basis independent, but $\Xi$ is not. The Banach space $\mathcal{B}%
_{1}\left(  \mathcal{H}^{N}\right)  $ can be expressed as the direct sums
\end{subequations}
\begin{equation}
\mathcal{B}_{1}\left(  \mathcal{H}^{N}\right)  =\bigoplus_{i=0}^{N}%
\mathcal{B}_{i}^{N}%
\end{equation}
where
\begin{align}
\mathcal{B}_{i}^{N}  &  =KerL_{N}^{i-1}\circleddash KerL_{N}^{i};\qquad1\leq
i\leq N\nonumber\\
\mathcal{B}_{0}^{N}  &  =\left[  KerL_{N}^{0}\right]  ^{C}%
\end{align}
and $\left[  KerL_{N}^{0}\right]  ^{C}$ is the canonical compliment of
$KerL_{N}^{0}$ and
\begin{equation}
\mathcal{B}_{1}\left(  \mathcal{H}^{N}\right)  =Ker\Xi\oplus\left[
Ker\Xi\right]  ^{C}=Ker\Xi_{OD}\oplus Ker\Xi_{D}\oplus\left[  Ker\Xi\right]
^{C}.
\end{equation}
$\lceil$Note that $L_{N}^{0}\left(  X\right)  =Tr\left\{  X\right\}  \rfloor$.
We have the following relationship between kernels
\begin{align}
KerL_{N}^{0}  &  \supset Ker\Xi\supset KerL_{N}^{1}\tag{*}\\
&  \Downarrow\nonumber\\
\left[  KerL_{N}^{0}\right]  ^{C}  &  \subset\left[  Ker\Xi\right]
^{C}\subset\left[  KerL_{N}^{1}\right]  ^{C}%
\end{align}

\end{enumerate}

\begin{itemize}
\item[*] One can define semi groups of $\mathfrak{S}$ that leave these kernels
invariant by
\begin{subequations}
\begin{align}
\mathfrak{S}_{\Xi}  &  =\left\{  \breve{T}|\Xi\left(  \breve{T}X\breve
{T}^{\dagger}\right)  =0\text{ if }\Xi\left(  X\right)  =0\right\} \\
\mathfrak{S}_{L_{N}^{1}}  &  =\left\{  \breve{T}|L_{N}^{1}\left(  \breve
{T}X\breve{T}^{\dagger}\right)  =0\text{ if }L_{N}^{1}\left(  X\right)
=0\right\} \\
\mathfrak{S}_{L_{N}^{2}}  &  =\left\{  \breve{T}|L_{N}^{2}\left(  \breve
{T}X\breve{T}^{\dagger}\right)  =0\text{ if }L_{N}^{2}\left(  X\right)
=0\right\}
\end{align}
Thus
\end{subequations}
\begin{subequations}
\begin{align}
\Xi\left(  \breve{T}X\breve{T}^{\dagger}\right)   &  =\Xi\left(  X\right)
\quad\forall\breve{T}\in\mathfrak{S}_{\Xi}\\
L_{N}^{1}\left(  \breve{T}X\breve{T}^{\dagger}\right)   &  =L_{N}^{1}\left(
X\right)  \quad\forall\breve{T}\in\mathfrak{S}_{L_{N}^{1}}\\
L_{N}^{2}\left(  \breve{T}X\breve{T}^{\dagger}\right)   &  =L_{N}^{2}\left(
X\right)  \quad\forall\breve{T}\in\mathfrak{S}_{L_{N}^{2}}%
\end{align}

\end{subequations}
\end{itemize}

\begin{enumerate}
Any operator that belongs to these subspaces (the kernels of the contraction
maps) with an integral kernel that is non zero almost everywhere can be used
as a reference operator to generate all other operators belonging to these
subspaces. The integral kernels of the semi groups associated with the
contraction maps are characterized by the following:
\end{enumerate}

\begin{itemize}
\item The density contraction: in order that $Ker\Xi\rightarrow Ker\Xi$ by
$\mathfrak{S}_{\Xi}$ we must have that
\begin{equation}
\int X\left(  \mathsf{z}^{N}\right)  \left\vert T\left(  \mathsf{z}%
^{N}\right)  \right\vert ^{2}d\mathsf{z}_{2}\cdots d\mathsf{z}_{N}%
\overset{a.e.}{=}0\qquad\forall\breve{X}\in Ker\Xi
\end{equation}
i.e.
\begin{equation}
\hat{T}:Ker\Xi_{D}\rightarrow Ker\Xi_{D}%
\end{equation}
$\lceil$Note that $\hat{T}:Ker\Xi_{OD}\rightarrow Ker\Xi_{OD}$ automatically
so this is not a constraint$\rfloor$.

\item The one density contraction: in order that $KerL_{N}^{1}\rightarrow
KerL_{N}^{1}$ by $\mathfrak{S}_{L_{N}^{1}}$ we must have that%
\begin{equation}
\int X\left(  \mathsf{z}_{1}\mathsf{z}^{N-1}\mathsf{,z}_{1}^{\prime}%
\mathsf{z}^{N-1}\right)  T\left(  \mathsf{z}_{1}\mathsf{z}^{N-1}\right)
T\left(  \mathsf{z}_{1}^{\prime}\mathsf{z}^{N-1}\right)  ^{\ast}%
d\mathsf{z}^{N-1}\overset{a.e.}{=}0\qquad\forall\breve{X}\in KerL_{N}^{1}%
\end{equation}


\item The two density contraction: in order that $KerL_{N}^{2}\rightarrow
KerL_{N}^{2}$ by $\mathfrak{S}_{L_{N}^{2}}$ we must have that%
\begin{equation}
\int X\left(  \mathsf{z}_{1}\mathsf{z}_{2}\mathsf{z}^{N-2}\mathsf{,z}%
_{1}^{\prime}\mathsf{z}_{2}^{\prime}\mathsf{z}^{N-2}\right)  T\left(
\mathsf{z}_{1}\mathsf{z}_{2}\mathsf{z}^{N-2}\right)  T\left(  \mathsf{z}%
_{1}^{\prime}\mathsf{z}_{2}^{\prime}\mathsf{z}^{N-2}\right)  ^{\ast
}d\mathsf{z}^{N-2}\overset{a.e.}{=}0\,\forall\breve{X}\in KerL_{N}^{2}%
\end{equation}


\item[*] One can also construct semi groups that have the property%
\begin{subequations}
\begin{align}
\mathfrak{S}_{\Xi}^{C}  &  :\left[  Ker\Xi\right]  ^{C}\rightarrow\left[
Ker\Xi\right]  ^{C}\\
\mathfrak{S}_{L_{N}^{1}}^{C}  &  :\left[  KerL_{N}^{1}\right]  ^{C}%
\rightarrow\left[  KerL_{N}^{1}\right]  ^{C}\\
\mathfrak{S}_{L_{N}^{2}}^{C}  &  :\left[  KerL_{N}^{2}\right]  ^{C}%
\rightarrow\left[  KerL_{N}^{2}\right]  ^{C}%
\end{align}
and any operator that belongs to these subspaces with an integral kernel that
is non zero almost everywhere can be used as a reference to generate all other
operators belonging to these subspaces.

\item[*] One can construct vector bundles with fibers that are the kernels of
the contraction maps and the base of the fibers being the canonical
compliments of the kernels.

\item[*] Using the groups $\left\{  \mathfrak{S}_{\mathcal{X}}^{C}%
|\mathcal{X=}\Xi,L_{N}^{1},L_{N}^{2}\right\}  $and some technical definitions
of coset type objects one can define maps that uniquely take a reference
operator in $\mathcal{B}_{1+}\left(  \mathcal{H}^{N}\right)  $ to the root of
each fiber, $T_{B}D_{ref}^{N}T_{B}^{\dagger}$, in an unique manner and then by
using cosets of $\left\{  \mathfrak{S}_{\mathcal{X}}|\mathcal{X=}\Xi,L_{N}%
^{1},L_{N}^{2}\right\}  $ uniquely from $T_{B}D_{ref}^{N}T_{B}^{\dagger}$ to
any state in the fiber $\left[  T_{B}D_{ref}^{N}T_{B}^{\dagger}\right]  _{N}$.

\item[*] The energy functional is defined on $\mathcal{B}_{1+}\left(
\mathcal{H}^{N}\right)  \rightarrow\mathbb{R}$ by
\end{subequations}
\begin{equation}
E\left(  D^{N}\right)  =\frac{Tr\left\{  HD^{N}\right\}  }{Tr\left\{
D^{N}\right\}  }%
\end{equation}
from which one can define energy functionals on the base as
\begin{equation}
\mathcal{E}\left(  T_{B}\right)  =\underset{D^{N}\in\left[  T_{B}D_{ref}%
^{N}T_{B}^{\dagger}\right]  _{N}}{Min}E\left(  D^{N}\right)  =\underset
{T_{F}\in\mathfrak{S}_{\mathcal{X}}}{Min}E\left(  T_{F}T_{B}D_{ref}^{N}%
T_{B}^{\dagger}T_{F}^{\dagger}\right)
\end{equation}
where $T_{B}\in\mathfrak{S}_{\mathcal{X}}^{C}$ and $T_{F}\in\mathfrak{S}%
_{\mathcal{X}}$ for $\mathcal{X=}\Xi$ or $L_{N}^{1}$ or $L_{N}^{2}$ and we
have used the non uniqueness of $Ker\Xi^{C}$ to represent the base in non
canonical form as $\left\{  T_{B}D_{ref}^{N}T_{B}^{\dagger}|T_{B}%
\in\mathfrak{S}_{\mathcal{X}}^{C}\right\}  $.

\item[*] In the time dependent case these would be action functionals.
\end{itemize}


\end{document} 